% GENERATED FILE. Edit canonical lessons, then run npm run generate:books.
% canonical-source-hash: fnv1a64:0969e125e96dcf78
% canonical-lessons: KA-C301-gaju, KA-C301-plastik, KA-C301-ittige, KA-C301-patti, KA-C301-kavali

\chapter{Glass, Plastic, Brick, List, Griddle}
\label{ch:ka-glass-plastic-brick-list-griddle}
\hlchaptermodality{301}

\begin{chapteropening}
\textbf{By the end of this chapter:} \emph{I can say glass, plastic, a brick, a list and a griddle.}

Complete the last lesson of chapter 301: I can say glass, plastic, a brick, a list and a griddle.
\end{chapteropening}

\section[\texorpdfstring{\kn{ಗಾಜು} (gāju)}{ಗಾಜು (gāju)}]{\kn{ಗಾಜು} (gāju) — glass (the material)}

\label{lesson:KA-C301-gaju}



\begin{quote}
Before the new one: say the Kannada for a TV serial, then the Kannada for a novel.
\end{quote}

\subsection*{You'll want to know: \kn{ಗಾಜು}}
\textbf{\kn{ಗಾಜು}} — \emph{gāju} — \textquotedblleft{}glass (the material)\textquotedblright{}.

\begin{grammarlens}[title={What you've built}]
One more word for this chapter.
\end{grammarlens}

\subsection*{Your turn}
\begin{itemize}
\raggedright
  \item \emph{Say it:} \emph{gāju}
  \item \emph{Say it:} \emph{gāju}, once more
  \item \emph{Say it:} say \emph{gāju}
  \item \emph{Recall:} say the Kannada for a TV serial, then the Kannada for a novel, then say \emph{gāju} again
\end{itemize}

\begin{tcolorbox}[breakable,colback=teal!4,colframe=teal!35!black,title={Before you move on}]
What does \kn{ಗಾಜು} mean? (\textquotedblleft{}Glass (the material)\textquotedblright{}.) Say it once more.
\end{tcolorbox}
\section[\texorpdfstring{\kn{ಪ್ಲಾಸ್ಟಿಕ್} (plāsṭik)}{ಪ್ಲಾಸ್ಟಿಕ್ (plāsṭik)}]{\kn{ಪ್ಲಾಸ್ಟಿಕ್} (plāsṭik) — plastic}

\label{lesson:KA-C301-plastik}



\begin{quote}
Before the new one: say the Kannada for a novel, then the Kannada for glass.
\end{quote}

\subsection*{You'll want to know: \kn{ಪ್ಲಾಸ್ಟಿಕ್}}
\textbf{\kn{ಪ್ಲಾಸ್ಟಿಕ್}} — \emph{plāsṭik} — \textquotedblleft{}plastic\textquotedblright{}.

\begin{grammarlens}[title={What you've built}]
One more word for this chapter.
\end{grammarlens}

\subsection*{Your turn}
\begin{itemize}
\raggedright
  \item \emph{Say it:} \emph{plāsṭik}
  \item \emph{Say it:} \emph{plāsṭik}, once more
  \item \emph{Say it:} say \emph{plāsṭik}
  \item \emph{Recall:} say the Kannada for a novel, then the Kannada for glass, then say \emph{plāsṭik} again
\end{itemize}

\begin{tcolorbox}[breakable,colback=teal!4,colframe=teal!35!black,title={Before you move on}]
What does \kn{ಪ್ಲಾಸ್ಟಿಕ್} mean? (\textquotedblleft{}Plastic\textquotedblright{}.) Say it once more.
\end{tcolorbox}
\section[\texorpdfstring{\kn{ಇಟ್ಟಿಗೆ} (iṭṭige)}{ಇಟ್ಟಿಗೆ (iṭṭige)}]{\kn{ಇಟ್ಟಿಗೆ} (iṭṭige) — a brick}

\label{lesson:KA-C301-ittige}



\begin{quote}
Before the new one: say the Kannada for glass, then the Kannada for plastic.
\end{quote}

\subsection*{You'll want to know: \kn{ಇಟ್ಟಿಗೆ}}
\textbf{\kn{ಇಟ್ಟಿಗೆ}} — \emph{iṭṭige} — \textquotedblleft{}a brick\textquotedblright{}.

\begin{grammarlens}[title={What you've built}]
One more word for this chapter.
\end{grammarlens}

\subsection*{Your turn}
\begin{itemize}
\raggedright
  \item \emph{Say it:} \emph{iṭṭige}
  \item \emph{Say it:} \emph{iṭṭige}, once more
  \item \emph{Say it:} say \emph{iṭṭige}
  \item \emph{Recall:} say the Kannada for glass, then the Kannada for plastic, then say \emph{iṭṭige} again
\end{itemize}

\begin{tcolorbox}[breakable,colback=teal!4,colframe=teal!35!black,title={Before you move on}]
What does \kn{ಇಟ್ಟಿಗೆ} mean? (\textquotedblleft{}A brick\textquotedblright{}.) Say it once more.
\end{tcolorbox}
\section[\texorpdfstring{\kn{ಪಟ್ಟಿ} (paṭṭi)}{ಪಟ್ಟಿ (paṭṭi)}]{\kn{ಪಟ್ಟಿ} (paṭṭi) — a list}

\label{lesson:KA-C301-patti}



\begin{quote}
Before the new one: say the Kannada for plastic, then the Kannada for a brick.
\end{quote}

\subsection*{You'll want to know: \kn{ಪಟ್ಟಿ}}
\textbf{\kn{ಪಟ್ಟಿ}} — \emph{paṭṭi} — \textquotedblleft{}a list\textquotedblright{}.

\begin{grammarlens}[title={What you've built}]
One more word for this chapter.
\end{grammarlens}

\subsection*{Your turn}
\begin{itemize}
\raggedright
  \item \emph{Say it:} \emph{paṭṭi}
  \item \emph{Say it:} \emph{paṭṭi}, once more
  \item \emph{Say it:} say \emph{paṭṭi}
  \item \emph{Recall:} say the Kannada for plastic, then the Kannada for a brick, then say \emph{paṭṭi} again
\end{itemize}

\begin{tcolorbox}[breakable,colback=teal!4,colframe=teal!35!black,title={Before you move on}]
What does \kn{ಪಟ್ಟಿ} mean? (\textquotedblleft{}A list\textquotedblright{}.) Say it once more.
\end{tcolorbox}
\section[\texorpdfstring{\kn{ಕಾವಲಿ} (kāvali)}{ಕಾವಲಿ (kāvali)}]{\kn{ಕಾವಲಿ} (kāvali) — a griddle, a dosa pan}

\label{lesson:KA-C301-kavali}



\begin{quote}
Before the new one: say the Kannada for a brick, then the Kannada for a list.
\end{quote}

\subsection*{You'll want to know: \kn{ಕಾವಲಿ}}
\textbf{\kn{ಕಾವಲಿ}} — \emph{kāvali} — \textquotedblleft{}a griddle, a dosa pan\textquotedblright{}.

\begin{grammarlens}[title={What you've built}]
One more word for this chapter.
\end{grammarlens}

\subsection*{Your turn}
\begin{itemize}
\raggedright
  \item \emph{Say it:} \emph{kāvali}
  \item \emph{Say it:} \emph{kāvali}, once more
  \item \emph{Say it:} say \emph{kāvali}
  \item \emph{Recall:} say the Kannada for a brick, then the Kannada for a list, then say \emph{kāvali} again
\end{itemize}

\begin{tcolorbox}[breakable,colback=teal!4,colframe=teal!35!black,title={Before you move on}]
What does \kn{ಕಾವಲಿ} mean? (\textquotedblleft{}A griddle, a dosa pan\textquotedblright{}.) Say it once more.
\end{tcolorbox}
